\section{Introduction}
\label{sec:introduction}

Modern cybersecurity monitoring produces large volumes of heterogeneous logs: firewall events, IDS alerts, application traces, HTTP requests, user-agent strings, actions, transferred bytes, and threat labels. A batch-only system can explain what happened historically, but it cannot alert quickly enough when an attack is active. A streaming-only system gives fresh alerts, but it does not provide the long-term reputation and historical context needed to prioritize decisions. For this reason, the project statement asks for a complete Lambda Architecture: batch processing, speed processing, serving integration, and dashboard visualization.

RAPID answers this need by building an end-to-end system around the provided cybersecurity log dataset. The same fields required by the assignment are used throughout the pipeline: \texttt{timestamp}, \texttt{source\_ip}, \texttt{dest\_ip}, \texttt{protocol}, \texttt{action}, \texttt{threat\_label}, \texttt{log\_type}, \texttt{bytes\_transferred}, \texttt{user\_agent}, and \texttt{request\_path}. These fields support both historical analytics and real-time rules: malicious IP ranking, port-scan detection, SQL injection and XSS signature extraction, brute-force detection, abnormal volume detection, and threat scoring.

\subsection{Objectives}

The project objectives are the following:

\begin{itemize}
    \item store the historical dataset in HDFS using a partitioned layout suitable for replay and batch analysis;
    \item compute historical attack patterns with Spark batch jobs;
    \item store batch views in HBase for fast serving;
    \item create a Kafka topic named \texttt{cybersecurity-logs} and simulate a continuous log stream;
    \item detect active threats with Spark Structured Streaming;
    \item store real-time threats in Cassandra;
    \item expose historical and active threats through a REST API;
    \item visualize the system through dashboards showing both historical and active threats.
\end{itemize}

\subsection{Requirement Coverage}

Table~\ref{tab:requirements-coverage} maps the assignment requirements to the RAPID implementation. This table is intentionally placed early because it gives the evaluator a direct checklist of the expected deliverables.

\begin{table}[H]
    \centering
    \scriptsize
    \caption{Coverage of the project statement requirements.}
    \label{tab:requirements-coverage}
    \begin{tabularx}{\textwidth}{>{\bfseries}p{0.22\textwidth} X X}
        \toprule
        Requirement & Expected by assignment & RAPID implementation \\
        \midrule
        Batch storage & HDFS folders partitioned by date & HDFS raw log partitions under \texttt{/logs/year=2024/month=XX/data.csv}; timestamps preserve day-level analysis and can be replayed by period. \\
        Batch analytics & Top malicious IPs, port scans, attack paths, volume by threat & Spark jobs compute IP reputation, port scans in 5-minute TCP windows, SQLi/XSS/path traversal/tool patterns, threat timeline, and bytes by threat label. \\
        Batch serving & HBase tables for reputation, patterns, timeline & HBase tables: \texttt{ip\_reputation}, \texttt{attack\_patterns}, \texttt{threat\_timeline}, \texttt{port\_scans}, \texttt{threat\_volume}, and \texttt{multistep\_attacks}. \\
        Speed ingestion & Kafka topic with 3 partitions and custom producer & Topic \texttt{cybersecurity-logs}; producer replays HDFS CSV partitions as JSON events using \texttt{source\_ip} as key. \\
        Real-time detection & Brute force, signatures, abnormal volume, score by IP & Spark Streaming jobs detect 5+ blocked attempts per minute, SQLMap/Nikto/Nmap/injection signatures, traffic-volume anomalies, and aggregate threat scores. \\
        Active threat store & Cassandra schema for recent threats & Cassandra keyspace \texttt{cybersecurity} with \texttt{logs}, \texttt{realtime\_threats}, \texttt{signature\_alerts}, \texttt{volume\_alerts}, and \texttt{threat\_scores}. \\
        Integration & API merging batch and speed results & Flask endpoint \texttt{/threats/ip/<ip>} combines HBase historical score and Cassandra active state, then returns score, level, and recommendation. \\
        Visualization & Dashboard for historical and active threats & HTML/JS dashboards show KPIs, timelines, IP reputation, port scans, attack patterns, active alerts, top risky sources, and a threat map. \\
        Bonus & ML, correlation, geolocation, adaptive scoring & Random Forest classification, multi-step attack correlation, geolocation-oriented map, and adaptive thresholds are documented and demonstrated. \\
        \bottomrule
    \end{tabularx}
\end{table}
